\section{Memory / I/O fabrics: CXL、PCIe IDE、RDMA}

\begin{frame}[t,shrink=6]{Memory / I/O fabrics: CXL、PCIe IDE、RDMA：方向开场}
\DeckKicker{方向开场}
\SlideMeta{12-memory-io-fabrics}{authored direction}{local evidence synthesis}
{\large\bfseries\color{Ink} 这个方向讨论的是高性能 memory/fabric data path 如何把 confidential boundary 推出 CPU 本地内存。\par}\vspace{1.8mm}
\begin{columns}[T,totalwidth=\textwidth]
  \begin{column}{0.45\textwidth}
    \fcolorbox{Rule}{Paper}{%
  \begin{minipage}[t]{0.97\linewidth}
    \vspace{1.1mm}
    \hspace{1.4mm}\begin{minipage}{0.92\linewidth}
      \PanelTitle{讲解逻辑}
      \scriptsize \begin{itemize}
  \item DirectCXL 证明 CXL.mem 可以让 host 通过 load/store 直接访问远端内存资源，核心变量是 latency 和 copy path。
  \item CXL-Tiers 讨论云虚拟化里本地 DRAM 与 CXL capacity tier 的 placement、隔离和尾部性能问题。
  \item ODRP 说明 RDMA remote paging 的瓶颈不只是数据传输，还包括远端内存分配、MNode CPU 与 utilization。
  \item 三篇都属于 data-path/fabric substrate：能解释为何需要跨 fabric 建模，但不能单独证明 CCA/CoVE confidential I/O 安全性。
\end{itemize}
    \end{minipage}
    \vspace{1mm}
  \end{minipage}}
  \end{column}
  \begin{column}{0.49\textwidth}
    \fcolorbox{Rule}{Panel}{%
  \begin{minipage}[t]{0.97\linewidth}
    \vspace{1.1mm}
    \hspace{1.4mm}\begin{minipage}{0.92\linewidth}
      \PanelTitle{三条 data-path 演进线}
      \scriptsize \textcolor{Accent}{\bfseries 2022 DirectCXL}\par\textbf{CXL.mem direct load/store}\par{\tiny\textcolor{Muted}{2705 cycles RDMA vs 328 cycles CXL}}\par\vspace{1.6mm}
\textcolor{Accent}{\bfseries 2024 CXL-Tiers}\par\textbf{CXL memory tiering for VMs}\par{\tiny\textcolor{Muted}{82\% workloads <=5\% slowdown}}\par\vspace{1.6mm}
\textcolor{Accent}{\bfseries 2025 ODRP}\par\textbf{programmable RDMA remote paging}\par{\tiny\textcolor{Muted}{1.72x-12x utilization gain}}\par\vspace{1.6mm}
    \end{minipage}
    \vspace{1mm}
  \end{minipage}}
  \end{column}
\end{columns}
\vspace{1mm}
{\tiny\textcolor{Muted}{Source: 本页只建立方向边界；性能数字分别来自三篇论文的摘要、Figure 6/10、Figure 5/10/12、Figure 7/9。}}
\end{frame}


\begin{frame}[t,shrink=6]{Direct Access, High-Performance Memory Disaggregation with DirectCXL：内容摘要}
\DeckKicker{内容摘要}
\SlideMeta{gouk2022directcxl}{E1 primary systems / Background substrate}{local\_pdf\_verified}
{\large\bfseries\color{Ink} DirectCXL 的核心贡献不是“更快网络”，而是把远端内存变成 host 可直接 load/store 的 CXL.mem 地址空间。\par}\vspace{1.8mm}
\begin{columns}[T,totalwidth=\textwidth]
  \begin{column}{0.45\textwidth}
    \fcolorbox{Rule}{Paper}{%
  \begin{minipage}[t]{0.97\linewidth}
    \vspace{1.1mm}
    \hspace{1.4mm}\begin{minipage}{0.92\linewidth}
      \PanelTitle{讲解逻辑}
      \scriptsize \begin{itemize}
  \item 论文提出 directly accessible memory disaggregation：host processor complex 通过 CXL.mem 直连 remote memory resources。
  \item 与 RDMA 路径相比，DirectCXL 不把数据先搬到 host DRAM 或 RNIC memory region，而是把 HDM 映射进系统地址空间。
  \item 作者还补了当时缺失的系统软件：Linux 5.13 runtime、cxl-namespace、segment table 和 mmap 接口。
  \item 在本 survey 中，它是 fabric data-path 基础论文，不是 confidential I/O 或 device trust 论文。
\end{itemize}
    \end{minipage}
    \vspace{1mm}
  \end{minipage}}
  \end{column}
  \begin{column}{0.49\textwidth}
    \fcolorbox{Rule}{Panel}{%
  \begin{minipage}[t]{0.97\linewidth}
    \vspace{1.1mm}
    \hspace{1.4mm}\begin{minipage}{0.92\linewidth}
      \PanelTitle{RDMA copy path vs CXL.mem direct path}
      \scriptsize \textbf{RDMA disaggregation}\par
page fault / object request → host runtime copies to RNIC MR → network transfer → memory-node runtime or RNIC → remote memory\par
{\tiny\textcolor{Muted}{copy + software fabric intervention}}\par\vspace{1.8mm}
\textbf{DirectCXL}\par
load/store → CXL root port → CXL switch → HDM address translation → remote DRAM\par
{\tiny\textcolor{Muted}{address-mapped access, no host-memory copy}}\par\vspace{1.8mm}
    \end{minipage}
    \vspace{1mm}
  \end{minipage}}
  \end{column}
\end{columns}
\vspace{1mm}
{\tiny\textcolor{Muted}{Source: 重绘自 DirectCXL Figure 1 与 Figure 2-Figure 4；未复制原论文图片。}}
\end{frame}


\begin{frame}[t,shrink=6]{Direct Access, High-Performance Memory Disaggregation with DirectCXL：研究背景}
\DeckKicker{研究背景}
\SlideMeta{gouk2022directcxl}{E1 primary systems / Background substrate}{local\_pdf\_verified}
{\large\bfseries\color{Ink} 传统 memory disaggregation 的主要成本来自 page/object 软件路径，而不是单纯的链路带宽。\par}\vspace{1.8mm}
\begin{columns}[T,totalwidth=\textwidth]
  \begin{column}{0.45\textwidth}
    \fcolorbox{Rule}{Paper}{%
  \begin{minipage}[t]{0.97\linewidth}
    \vspace{1.1mm}
    \hspace{1.4mm}\begin{minipage}{0.92\linewidth}
      \PanelTitle{讲解逻辑}
      \scriptsize \begin{itemize}
  \item page-based 方案依赖虚拟内存和 swap/page-fault，把页从 host 与 memory node 之间搬运，透明但有 page fault、I/O amplification 和 context switch。
  \item object-based 方案用 key-value 或对象接口绕开部分地址转换成本，但要求应用改接口，语义依赖应用。
  \item 两类方案通常仍依赖 RDMA 或类似 DMA 协议，仍需要 memory region、metadata、copy 和 runtime 参与。
  \item 因此论文的背景问题是：能否让远端内存像系统内存一样被 CPU 直接访问。
\end{itemize}
    \end{minipage}
    \vspace{1mm}
  \end{minipage}}
  \end{column}
  \begin{column}{0.49\textwidth}
    \fcolorbox{Rule}{Panel}{%
  \begin{minipage}[t]{0.97\linewidth}
    \vspace{1.1mm}
    \hspace{1.4mm}\begin{minipage}{0.92\linewidth}
      \PanelTitle{传统路径的三类开销}
      \scriptsize {\tiny
\begin{tabularx}{\linewidth}{@{}YYYY@{}}
\toprule
\textbf{路径} & \textbf{优势} & \textbf{主要成本} & \textbf{对 confidential data path 的含义} \\
\midrule
Page-based & 应用透明 & page fault、swap、copy、context switch & 敏感页会被软件路径和远端节点状态影响 \\
Object-based & 可避开部分 VM 开销 & 应用改造、对象语义绑定 & 只能保护被重写接口覆盖的数据 \\
DirectCXL 目标 & load/store-like & CXL/PCIe latency 和 HDM 管理 & fabric 成为内存边界的一部分 \\
\bottomrule
\end{tabularx}
}
    \end{minipage}
    \vspace{1mm}
  \end{minipage}}
  \end{column}
\end{columns}
\vspace{1mm}
{\tiny\textcolor{Muted}{Source: 矩阵依据 DirectCXL Section 1-2 的 related-work 分类和 Figure 1 的 RDMA 路径重绘。}}
\end{frame}


\begin{frame}[t,shrink=6]{Direct Access, High-Performance Memory Disaggregation with DirectCXL：解决方案}
\DeckKicker{解决方案}
\SlideMeta{gouk2022directcxl}{E1 primary systems / Background substrate}{local\_pdf\_verified}
{\large\bfseries\color{Ink} DirectCXL 的设计妙处是把远端 DRAM 暴露为 HDM，并用 CXL switch/runtime 管理 namespace，而不是让远端 CPU 管数据。\par}\vspace{1.8mm}
\begin{columns}[T,totalwidth=\textwidth]
  \begin{column}{0.45\textwidth}
    \fcolorbox{Rule}{Paper}{%
  \begin{minipage}[t]{0.97\linewidth}
    \vspace{1.1mm}
    \hspace{1.4mm}\begin{minipage}{0.92\linewidth}
      \PanelTitle{讲解逻辑}
      \scriptsize \begin{itemize}
  \item CXL device 是被动内存模块：CXL controller 解析 CXL flits，将地址转换后交给内部 DRAM controllers。
  \item host kernel 在 PCIe enumeration 后映射 BAR/HDM，CPU 访问 HDM system memory 时请求经 root port 转成 CXL flit。
  \item CXL switch 用 virtual hierarchy 和 routing table 连接 host 与 CXL device；runtime 用 /dev/directcxl、ioctl、cxl-namespace 和 mmap 暴露给应用。
  \item 这个方案的逻辑是把 copy path 改成 address path，把 memory node computing resource 从关键路径拿掉。
\end{itemize}
    \end{minipage}
    \vspace{1mm}
  \end{minipage}}
  \end{column}
  \begin{column}{0.49\textwidth}
    \fcolorbox{Rule}{Panel}{%
  \begin{minipage}[t]{0.97\linewidth}
    \vspace{1.1mm}
    \hspace{1.4mm}\begin{minipage}{0.92\linewidth}
      \PanelTitle{DirectCXL 访问机制}
      \scriptsize \begin{itemize}
  \item 01. 应用 mmap cxl-namespace
  \item 02. CPU load/store 到 HDM system memory
  \item 03. Root Port 转成 CXL flit
  \item 04. CXL switch virtual hierarchy 路由
  \item 05. CXL controller 访问远端 DRAM
  \item 06. 数据沿 CXL path 返回 CPU cache
\end{itemize}
    \end{minipage}
    \vspace{1mm}
  \end{minipage}}
  \end{column}
\end{columns}
\vspace{1mm}
{\tiny\textcolor{Muted}{Source: 机制图按 DirectCXL Figure 2-Figure 5 重绘，保留 HDM、switch、runtime 三个关键对象。}}
\end{frame}


\begin{frame}[t,shrink=6]{Direct Access, High-Performance Memory Disaggregation with DirectCXL：实验结果}
\DeckKicker{实验结果}
\SlideMeta{gouk2022directcxl}{E1 primary systems / Background substrate}{local\_pdf\_verified}
{\large\bfseries\color{Ink} DirectCXL 的实测收益集中在 64B load latency 与真实 workload：RDMA 2705 cycles，DirectCXL 328 cycles，并带来约 3x 应用性能。\par}\vspace{1.8mm}
\begin{columns}[T,totalwidth=\textwidth]
  \begin{column}{0.45\textwidth}
    \fcolorbox{Rule}{Paper}{%
  \begin{minipage}[t]{0.97\linewidth}
    \vspace{1.1mm}
    \hspace{1.4mm}\begin{minipage}{0.92\linewidth}
      \PanelTitle{讲解逻辑}
      \scriptsize \begin{itemize}
  \item 64B read latency 中，RDMA 路径约 2705 cycles，DirectCXL memory load request 约 328 cycles，论文报告 8.3x faster。
  \item 从 memory hierarchy 看，RDMA 比 DirectCXL 慢 6.2x；DirectCXL 仍比 local DRAM 慢，瓶颈主要来自 PCIe/CXL IP。
  \item 真实 workload 中，DirectCXL 相对 RDMA-based memory disaggregation 平均约 3x 性能提升。
  \item 这些数据证明的是 fabric access path 的性能差异，不证明链路加密、设备身份或多租户隔离。
\end{itemize}
    \end{minipage}
    \vspace{1mm}
  \end{minipage}}
  \end{column}
  \begin{column}{0.49\textwidth}
    \fcolorbox{Rule}{Panel}{%
  \begin{minipage}[t]{0.97\linewidth}
    \vspace{1.1mm}
    \hspace{1.4mm}\begin{minipage}{0.92\linewidth}
      \PanelTitle{DirectCXL 关键实验数字}
      \scriptsize \textbf{RDMA 64B read}\hfill \textcolor{Accent}{\bfseries 2705 cycles}\par\textcolor{Accent}{\rule{0.95\linewidth}{1.1mm}}\par{\tiny\textcolor{Muted}{baseline}}\par\vspace{1.4mm}
\textbf{DirectCXL 64B load}\hfill \textcolor{Accent}{\bfseries 328 cycles}\par\textcolor{Accent}{\rule{0.12\linewidth}{1.1mm}}\par{\tiny\textcolor{Muted}{8.3x faster in Fig.6}}\par\vspace{1.4mm}
\textbf{RDMA vs DirectCXL hierarchy}\hfill \textcolor{Accent}{\bfseries 6.2x slower}\par\textcolor{Accent}{\rule{0.62\linewidth}{1.1mm}}\par{\tiny\textcolor{Muted}{Fig.8/Table 2}}\par\vspace{1.4mm}
\textbf{real workload speedup}\hfill \textcolor{Accent}{\bfseries about 3x}\par\textcolor{Accent}{\rule{0.45\linewidth}{1.1mm}}\par{\tiny\textcolor{Muted}{Fig.10/conclusion}}\par\vspace{1.4mm}
    \end{minipage}
    \vspace{1mm}
  \end{minipage}}
  \end{column}
\end{columns}
\vspace{1mm}
{\tiny\textcolor{Muted}{Source: 实验图按 DirectCXL Figure 6、Figure 8/Table 2、Figure 10 的数字重绘。}}
\end{frame}


\begin{frame}[t,shrink=6]{Direct Access, High-Performance Memory Disaggregation with DirectCXL：文章评价}
\DeckKicker{文章评价}
\SlideMeta{gouk2022directcxl}{E1 primary systems / Background substrate}{local\_pdf\_verified}
{\large\bfseries\color{Ink} DirectCXL 很适合说明 CXL 会改变内存边界，但它缺少 confidential I/O 所需的身份、加密和 lifecycle 证据。\par}\vspace{1.8mm}
\begin{columns}[T,totalwidth=\textwidth]
  \begin{column}{0.45\textwidth}
    \fcolorbox{Rule}{Paper}{%
  \begin{minipage}[t]{0.97\linewidth}
    \vspace{1.1mm}
    \hspace{1.4mm}\begin{minipage}{0.92\linewidth}
      \PanelTitle{讲解逻辑}
      \scriptsize \begin{itemize}
  \item 优点：问题切得准，系统原型覆盖 CXL device、switch、runtime 和 workload 评估，能支撑 memory fabric substrate 的叙事。
  \item 不足：论文目标是性能和 disaggregation，不覆盖 CXL IDE、SPDM/TDISP、attestation、multi-tenant memory ownership 或 verifier policy。
  \item 商业潜力：CXL memory expansion/pooling 有明确数据中心价值，DirectCXL 的 namespace/runtime 思路可转化为平台管理接口。
  \item 落地风险：真实商用环境还需要链路保护、设备身份、热插拔/故障恢复、资源隔离和运维可观测性。
\end{itemize}
    \end{minipage}
    \vspace{1mm}
  \end{minipage}}
  \end{column}
  \begin{column}{0.49\textwidth}
    \fcolorbox{Rule}{Panel}{%
  \begin{minipage}[t]{0.97\linewidth}
    \vspace{1.1mm}
    \hspace{1.4mm}\begin{minipage}{0.92\linewidth}
      \PanelTitle{商业化判断}
      \scriptsize {\tiny
\begin{tabularx}{\linewidth}{@{}YY@{}}
\toprule
\textbf{维度} & \textbf{判断} \\
\midrule
强项 & CXL.mem direct access 和 runtime 原型可解释 memory pooling 性能价值 \\
短板 & 不处理 device identity、link encryption、TDISP 或 verifier policy \\
可落地场景 & 云内存扩展、CXL pool 管理、远端内存性能建模 \\
不能承担的 claim & 不能单独证明 confidential I/O 或 CCA/CoVE data-path 安全 \\
\bottomrule
\end{tabularx}
}
    \end{minipage}
    \vspace{1mm}
  \end{minipage}}
  \end{column}
\end{columns}
\vspace{1mm}
{\tiny\textcolor{Muted}{Source: 评价结合 DirectCXL conclusion 与本仓库 evidence boundary，不添加论文未声称的安全结论。}}
\end{frame}


\begin{frame}[t,shrink=6]{Managing Memory Tiers with CXL in Virtualized Environments：内容摘要}
\DeckKicker{内容摘要}
\SlideMeta{zhong2024cxltiers}{E1 primary systems / Background substrate}{local\_pdf\_verified}
{\large\bfseries\color{Ink} CXL-Tiers 的关键不是“用 CXL 扩内存”，而是在 VM 云环境里把硬件 tiering 与软件隔离策略组合起来。\par}\vspace{1.8mm}
\begin{columns}[T,totalwidth=\textwidth]
  \begin{column}{0.45\textwidth}
    \fcolorbox{Rule}{Paper}{%
  \begin{minipage}[t]{0.97\linewidth}
    \vspace{1.1mm}
    \hspace{1.4mm}\begin{minipage}{0.92\linewidth}
      \PanelTitle{讲解逻辑}
      \scriptsize \begin{itemize}
  \item 论文提出 Intel Flat Memory Mode：在 memory controller 内以 cache-line 粒度管理 local DRAM 与 CXL memory。
  \item 硬件 tiering 避免 hypervisor/OS 用 page-table 扫描、PTE poisoning 或 instruction sampling 追踪 hotness。
  \item Memstrata 是轻量 multi-tenant allocator，用 slowdown estimator 和 dedicated local pages 缓解 outlier VM。
  \item 在本 survey 中，它说明 fabric memory 与 VM resource management 的冲突，但仍不是 confidential boundary 证明。
\end{itemize}
    \end{minipage}
    \vspace{1mm}
  \end{minipage}}
  \end{column}
  \begin{column}{0.49\textwidth}
    \fcolorbox{Rule}{Panel}{%
  \begin{minipage}[t]{0.97\linewidth}
    \vspace{1.1mm}
    \hspace{1.4mm}\begin{minipage}{0.92\linewidth}
      \PanelTitle{硬件 tiering + 软件 isolation}
      \scriptsize \begin{itemize}
  \item 01. CXL capacity tier 扩展 VM 可见内存
  \item 02. Memory controller cache-line swap local/CXL lines
  \item 03. mixed mode 保留 dedicated local DRAM
  \item 04. Memstrata 预测 VM slowdown
  \item 05. allocator 把 local pages 给 outlier VM
\end{itemize}
    \end{minipage}
    \vspace{1mm}
  \end{minipage}}
  \end{column}
\end{columns}
\vspace{1mm}
{\tiny\textcolor{Muted}{Source: 机制图重绘自 CXL-Tiers Figure 3、Figure 4 和 Figure 8。}}
\end{frame}


\begin{frame}[t,shrink=6]{Managing Memory Tiers with CXL in Virtualized Environments：研究背景}
\DeckKicker{研究背景}
\SlideMeta{zhong2024cxltiers}{E1 primary systems / Background substrate}{local\_pdf\_verified}
{\large\bfseries\color{Ink} 虚拟化让软件 memory tiering 很难做对：guest 不配合、page 粒度太粗、host 追踪成本又高。\par}\vspace{1.8mm}
\begin{columns}[T,totalwidth=\textwidth]
  \begin{column}{0.45\textwidth}
    \fcolorbox{Rule}{Paper}{%
  \begin{minipage}[t]{0.97\linewidth}
    \vspace{1.1mm}
    \hspace{1.4mm}\begin{minipage}{0.92\linewidth}
      \PanelTitle{讲解逻辑}
      \scriptsize \begin{itemize}
  \item 云平台希望用 CXL memory 降低成本和 embodied carbon，但 CXL 访问延迟高于本地 DRAM。
  \item 软件 tiering 需要追踪 hotness：PTE 扫描、PTE poisoning 或 instruction sampling 在 VM 场景里有开销、兼容性和隐私问题。
  \item page 级 placement 假设页内访问有空间局部性，但 VM 的 free-page randomization、warm VM 和 huge page 会破坏这个假设。
  \item 论文给出的设计目标包括 unmodified VM、低 host overhead、不增加 cross-VM interference、接近 local DRAM 性能。
\end{itemize}
    \end{minipage}
    \vspace{1mm}
  \end{minipage}}
  \end{column}
  \begin{column}{0.49\textwidth}
    \fcolorbox{Rule}{Panel}{%
  \begin{minipage}[t]{0.97\linewidth}
    \vspace{1.1mm}
    \hspace{1.4mm}\begin{minipage}{0.92\linewidth}
      \PanelTitle{VM 场景下 software tiering 的失配}
      \scriptsize {\tiny
\begin{tabularx}{\linewidth}{@{}YYY@{}}
\toprule
\textbf{问题} & \textbf{论文中的原因} & \textbf{对 confidential VM 的启发} \\
\midrule
host CPU 成本 & 细粒度 hotness 追踪要扫描 PTE 或触发 fault & 不能让不可信 host 为性能频繁窥视 guest 行为 \\
隐私/兼容性 & instruction sampling 不支持或不适合 VM & telemetry 本身可能成为信任边界问题 \\
placement 粒度 & 2MB/1GB page 会混合 hot/cold cache lines & ownership 与 placement 不能只按页粗糙绑定 \\
cross-VM interference & VM 会竞争 local DRAM capacity & 资源调度会影响机密 workload 尾部性能 \\
\bottomrule
\end{tabularx}
}
    \end{minipage}
    \vspace{1mm}
  \end{minipage}}
  \end{column}
\end{columns}
\vspace{1mm}
{\tiny\textcolor{Muted}{Source: 背景矩阵依据 CXL-Tiers Section 2 与 Figure 1-Figure 2 重绘。}}
\end{frame}


\begin{frame}[t,shrink=6]{Managing Memory Tiers with CXL in Virtualized Environments：解决方案}
\DeckKicker{解决方案}
\SlideMeta{zhong2024cxltiers}{E1 primary systems / Background substrate}{local\_pdf\_verified}
{\large\bfseries\color{Ink} Flat Memory Mode 把细粒度 placement 放到 memory controller，Memstrata 只处理 VM 间 performance isolation。\par}\vspace{1.8mm}
\begin{columns}[T,totalwidth=\textwidth]
  \begin{column}{0.45\textwidth}
    \fcolorbox{Rule}{Paper}{%
  \begin{minipage}[t]{0.97\linewidth}
    \vspace{1.1mm}
    \hspace{1.4mm}\begin{minipage}{0.92\linewidth}
      \PanelTitle{讲解逻辑}
      \scriptsize \begin{itemize}
  \item hardware-tiered NUMA node 将 local DRAM 与 CXL memory 组合成软件可见的 aggregate capacity，并在 cache-line 粒度交换。
  \item mixed mode 额外暴露 dedicated local memory NUMA node，为 cache-unfriendly 或 latency-sensitive VM 预留稳定本地页。
  \item Memstrata 读取 performance events，用 random forest slowdown estimator 识别 outlier VM，并迁移 dedicated local pages。
  \item 妙处是分工清楚：硬件处理行级 hotness，软件处理多租户公平和 outlier 缓解。
\end{itemize}
    \end{minipage}
    \vspace{1mm}
  \end{minipage}}
  \end{column}
  \begin{column}{0.49\textwidth}
    \fcolorbox{Rule}{Panel}{%
  \begin{minipage}[t]{0.97\linewidth}
    \vspace{1.1mm}
    \hspace{1.4mm}\begin{minipage}{0.92\linewidth}
      \PanelTitle{Memstrata 控制回路}
      \scriptsize \begin{itemize}
  \item 01. VM runs on HW-tiered + dedicated local memory
  \item 02. collect per-VM performance events
  \item 03. predict slowdown >5\%
  \item 04. rank outlier VMs
  \item 05. migrate dedicated local pages
  \item 06. reduce tail slowdown without guest cooperation
\end{itemize}
    \end{minipage}
    \vspace{1mm}
  \end{minipage}}
  \end{column}
\end{columns}
\vspace{1mm}
{\tiny\textcolor{Muted}{Source: 控制回路按 CXL-Tiers Figure 8、Figure 9 和 Listing 1 重绘。}}
\end{frame}


\begin{frame}[t,shrink=6]{Managing Memory Tiers with CXL in Virtualized Environments：实验结果}
\DeckKicker{实验结果}
\SlideMeta{zhong2024cxltiers}{E1 primary systems / Background substrate}{local\_pdf\_verified}
{\large\bfseries\color{Ink} CXL-Tiers 的实验结论很具体：mixed mode 中 82\% workloads 不超过 5\% slowdown，Memstrata 把 realistic multi-VM outlier 从 35\% 压到 <6\%。\par}\vspace{1.8mm}
\begin{columns}[T,totalwidth=\textwidth]
  \begin{column}{0.45\textwidth}
    \fcolorbox{Rule}{Paper}{%
  \begin{minipage}[t]{0.97\linewidth}
    \vspace{1.1mm}
    \hspace{1.4mm}\begin{minipage}{0.92\linewidth}
      \PanelTitle{讲解逻辑}
      \scriptsize \begin{itemize}
  \item 单 VM 115 workloads 中，mixed mode 下 82\% workload slowdown <=5\%，95\% workload slowdown <=10\%，但 outlier 可达 34\%。
  \item multi-VM 场景中，Memstrata 可以识别 outlier 并迁移 dedicated local pages，把 worst-case slowdown 从 35\% 降到小于 6\%。
  \item Memstrata 的最大 CPU overhead 是 single core 的 4\%，内存开销约 110 MB。
  \item 这些数据证明 CXL tiering 能接近 local DRAM，但也显示 resource-management policy 会决定尾部体验。
\end{itemize}
    \end{minipage}
    \vspace{1mm}
  \end{minipage}}
  \end{column}
  \begin{column}{0.49\textwidth}
    \fcolorbox{Rule}{Panel}{%
  \begin{minipage}[t]{0.97\linewidth}
    \vspace{1.1mm}
    \hspace{1.4mm}\begin{minipage}{0.92\linewidth}
      \PanelTitle{CXL-Tiers 关键实验数字}
      \scriptsize \textbf{mixed-mode workloads <=5\% slowdown}\hfill \textcolor{Accent}{\bfseries 82\%}\par\textcolor{Accent}{\rule{0.82\linewidth}{1.1mm}}\par{\tiny\textcolor{Muted}{115 workloads}}\par\vspace{1.4mm}
\textbf{mixed-mode workloads <=10\% slowdown}\hfill \textcolor{Accent}{\bfseries 95\%}\par\textcolor{Accent}{\rule{0.95\linewidth}{1.1mm}}\par{\tiny\textcolor{Muted}{single VM distribution}}\par\vspace{1.4mm}
\textbf{outlier before policy}\hfill \textcolor{Accent}{\bfseries up to 34-35\%}\par\textcolor{Accent}{\rule{0.35\linewidth}{1.1mm}}\par{\tiny\textcolor{Muted}{tail risk}}\par\vspace{1.4mm}
\textbf{after Memstrata}\hfill \textcolor{Accent}{\bfseries <6\%}\par\textcolor{Accent}{\rule{0.06\linewidth}{1.1mm}}\par{\tiny\textcolor{Muted}{realistic multi-VM}}\par\vspace{1.4mm}
\textbf{CPU overhead}\hfill \textcolor{Accent}{\bfseries 4\% of one core}\par\textcolor{Accent}{\rule{0.04\linewidth}{1.1mm}}\par{\tiny\textcolor{Muted}{manager + ML model}}\par\vspace{1.4mm}
    \end{minipage}
    \vspace{1mm}
  \end{minipage}}
  \end{column}
\end{columns}
\vspace{1mm}
{\tiny\textcolor{Muted}{Source: 实验条形图按 CXL-Tiers abstract、Figure 5、Figure 10-Figure 12 和正文数值重绘。}}
\end{frame}


\begin{frame}[t,shrink=6]{Managing Memory Tiers with CXL in Virtualized Environments：文章评价}
\DeckKicker{文章评价}
\SlideMeta{zhong2024cxltiers}{E1 primary systems / Background substrate}{local\_pdf\_verified}
{\large\bfseries\color{Ink} CXL-Tiers 最接近云平台落地，但它解决的是 performance isolation，不是机密内存 ownership 或 link trust。\par}\vspace{1.8mm}
\begin{columns}[T,totalwidth=\textwidth]
  \begin{column}{0.45\textwidth}
    \fcolorbox{Rule}{Paper}{%
  \begin{minipage}[t]{0.97\linewidth}
    \vspace{1.1mm}
    \hspace{1.4mm}\begin{minipage}{0.92\linewidth}
      \PanelTitle{讲解逻辑}
      \scriptsize \begin{itemize}
  \item 优点：真实 CXL hardware prototype、115 workloads、multi-VM outlier 实验和 Azure production motivation，让商业可信度强。
  \item 不足：Intel Flat Memory Mode 依赖特定处理器/MC 能力，且论文不解决 CXL IDE keying、attestation、device assignment 或 confidential VM ownership。
  \item 商业潜力：CXL capacity tier 能直接服务云成本、内存扩容和碳排优化，Memstrata 可作为 hypervisor 管理层策略。
  \item 对机密计算的价值：它提示 host placement policy 与 confidential boundary 之间会有冲突，需要把资源管理和安全证据分开写。
\end{itemize}
    \end{minipage}
    \vspace{1mm}
  \end{minipage}}
  \end{column}
  \begin{column}{0.49\textwidth}
    \fcolorbox{Rule}{Panel}{%
  \begin{minipage}[t]{0.97\linewidth}
    \vspace{1.1mm}
    \hspace{1.4mm}\begin{minipage}{0.92\linewidth}
      \PanelTitle{落地潜力与边界}
      \scriptsize {\tiny
\begin{tabularx}{\linewidth}{@{}YY@{}}
\toprule
\textbf{维度} & \textbf{判断} \\
\midrule
最强证据 & 真实 CXL prototype + 115 workloads + multi-VM 实验 \\
商业场景 & 云 VM memory expansion、DRAM cost/carbon reduction、tail-performance control \\
安全缺口 & 不覆盖 IDE/SPDM/TDISP、attestation、ownership lifecycle \\
对本 deck 的作用 & 支撑 fabric memory resource-management 问题，不支撑 confidential I/O 安全证明 \\
\bottomrule
\end{tabularx}
}
    \end{minipage}
    \vspace{1mm}
  \end{minipage}}
  \end{column}
\end{columns}
\vspace{1mm}
{\tiny\textcolor{Muted}{Source: 评价基于 CXL-Tiers conclusion/discussion 和本仓库 background-substrate 证据边界。}}
\end{frame}


\begin{frame}[t,shrink=6]{ODRP: on-demand remote paging with programmable RDMA：内容摘要}
\DeckKicker{内容摘要}
\SlideMeta{wang2025odrp}{E1 primary systems / Background substrate}{local\_pdf\_verified}
{\large\bfseries\color{Ink} ODRP 的贡献是把 remote paging 的内存管理也搬进 RNIC data path，而不是只用 RDMA 加速数据搬运。\par}\vspace{1.8mm}
\begin{columns}[T,totalwidth=\textwidth]
  \begin{column}{0.45\textwidth}
    \fcolorbox{Rule}{Paper}{%
  \begin{minipage}[t]{0.97\linewidth}
    \vspace{1.1mm}
    \hspace{1.4mm}\begin{minipage}{0.92\linewidth}
      \PanelTitle{讲解逻辑}
      \scriptsize \begin{itemize}
  \item 论文目标是 swap-based memory disaggregation：CNode 透明地把 cold pages swap 到 MNode。
  \item 传统 one-sided RDMA 适合读写远端 buffer，但不擅长动态 memory allocation/deallocation。
  \item ODRP 用 RDMA WR chain 在 RNIC 上实现 load、mapped store、unmapped store、invalidate 等 frontswap API。
  \item 结果是 page granularity remote memory management、zero MNode CPU usage 和接近 coarse-grained one-sided RDMA 的性能。
\end{itemize}
    \end{minipage}
    \vspace{1mm}
  \end{minipage}}
  \end{column}
  \begin{column}{0.49\textwidth}
    \fcolorbox{Rule}{Panel}{%
  \begin{minipage}[t]{0.97\linewidth}
    \vspace{1.1mm}
    \hspace{1.4mm}\begin{minipage}{0.92\linewidth}
      \PanelTitle{ODRP high-level path}
      \scriptsize \begin{itemize}
  \item 01. Linux frontswap request on CNode
  \item 02. CNode backend computes TT entry / request type
  \item 03. RDMA SEND triggers WR chain
  \item 04. RNIC reads TT and page queue
  \item 05. RNIC loads/stores/allocates 4KB page
  \item 06. MNode CPU stays off the data path
\end{itemize}
    \end{minipage}
    \vspace{1mm}
  \end{minipage}}
  \end{column}
\end{columns}
\vspace{1mm}
{\tiny\textcolor{Muted}{Source: 机制图按 ODRP Figure 3 和 Table 2 重绘。}}
\end{frame}


\begin{frame}[t,shrink=6]{ODRP: on-demand remote paging with programmable RDMA：研究背景}
\DeckKicker{研究背景}
\SlideMeta{wang2025odrp}{E1 primary systems / Background substrate}{local\_pdf\_verified}
{\large\bfseries\color{Ink} RDMA remote paging 的核心 tradeoff 是：static one-sided 快但浪费内存，dynamic/two-sided 节省内存却吃掉 MNode CPU。\par}\vspace{1.8mm}
\begin{columns}[T,totalwidth=\textwidth]
  \begin{column}{0.45\textwidth}
    \fcolorbox{Rule}{Paper}{%
  \begin{minipage}[t]{0.97\linewidth}
    \vspace{1.1mm}
    \hspace{1.4mm}\begin{minipage}{0.92\linewidth}
      \PanelTitle{讲解逻辑}
      \scriptsize \begin{itemize}
  \item one-sided static 预注册大 slab，runtime 不用 MNode CPU，但 internal fragmentation 导致 remote memory utilization 低。
  \item one-sided dynamic 细化 allocation 粒度，但 memory registration 是 control path 操作，4KB 或 1MB MR 注册代价很高。
  \item two-sided 可做到 4KB 粒度，但每次 swap 请求都要 MNode CPU 参与，CPU 很快成为瓶颈。
  \item ODRP 的研究问题是：能否同时获得 fine-grained utilization、zero MNode CPU 和 good performance。
\end{itemize}
    \end{minipage}
    \vspace{1mm}
  \end{minipage}}
  \end{column}
  \begin{column}{0.49\textwidth}
    \fcolorbox{Rule}{Panel}{%
  \begin{minipage}[t]{0.97\linewidth}
    \vspace{1.1mm}
    \hspace{1.4mm}\begin{minipage}{0.92\linewidth}
      \PanelTitle{RDMA alternatives tradeoff}
      \scriptsize {\tiny
\begin{tabularx}{\linewidth}{@{}YYYY@{}}
\toprule
\textbf{方案} & \textbf{memory utilization} & \textbf{MNode CPU} & \textbf{性能风险} \\
\midrule
one-sided static & 低，slab over-provision & 低 & 浪费 remote memory \\
one-sided dynamic & 中，1MB slab 仍有碎片 & 高，MR registration & allocation 队列和 CPU bottleneck \\
two-sided & 高，4KB 粒度 & 很高，每次 swap 参与 & CPU 饱和后延迟放大 \\
ODRP & 高，4KB page & 接近 0 & WR chain complexity \\
\bottomrule
\end{tabularx}
}
    \end{minipage}
    \vspace{1mm}
  \end{minipage}}
  \end{column}
\end{columns}
\vspace{1mm}
{\tiny\textcolor{Muted}{Source: 矩阵重绘自 ODRP Table 1、Figure 1 和 Section 2.2-2.3。}}
\end{frame}


\begin{frame}[t,shrink=6]{ODRP: on-demand remote paging with programmable RDMA：解决方案}
\DeckKicker{解决方案}
\SlideMeta{wang2025odrp}{E1 primary systems / Background substrate}{local\_pdf\_verified}
{\large\bfseries\color{Ink} ODRP 的设计妙处是 client-assisted WR chains：CNode 做轻量判断，RNIC 执行 TT/page-queue 操作，MNode CPU 不进关键路径。\par}\vspace{1.8mm}
\begin{columns}[T,totalwidth=\textwidth]
  \begin{column}{0.45\textwidth}
    \fcolorbox{Rule}{Paper}{%
  \begin{minipage}[t]{0.97\linewidth}
    \vspace{1.1mm}
    \hspace{1.4mm}\begin{minipage}{0.92\linewidth}
      \PanelTitle{讲解逻辑}
      \scriptsize \begin{itemize}
  \item MNode 把内存切成 4KB pages，用 FIFO free page queue 管空闲页，并为每个 CNode 维护 translation table。
  \item CNode backend 维护 allocation bitmap，提前判断 mapped/unmapped store，避免 WR chain 内额外分支。
  \item RNIC WR chain 用 READ/WRITE/CAS/FAA 等 primitives 实现 page load、mapped store、unmapped store、invalidate。
  \item EndianSwap、Modulo WR 和 WR recycling 是为克服 RNIC semantics 限制而做的工程设计。
\end{itemize}
    \end{minipage}
    \vspace{1mm}
  \end{minipage}}
  \end{column}
  \begin{column}{0.49\textwidth}
    \fcolorbox{Rule}{Panel}{%
  \begin{minipage}[t]{0.97\linewidth}
    \vspace{1.1mm}
    \hspace{1.4mm}\begin{minipage}{0.92\linewidth}
      \PanelTitle{CNode-assisted WR chain}
      \scriptsize \begin{itemize}
  \item 01. CNode computes tt\_addr and store type
  \item 02. RECV fills WR arguments
  \item 03. READ TT entry or free-page queue
  \item 04. FAA updates queue head/tail when allocating/freeing
  \item 05. WRITE page data and completion signal
  \item 06. CNode helps recycle consumed WRs
\end{itemize}
    \end{minipage}
    \vspace{1mm}
  \end{minipage}}
  \end{column}
\end{columns}
\vspace{1mm}
{\tiny\textcolor{Muted}{Source: WR-chain 图按 ODRP Figure 4-Figure 6 和 Section 4.3-4.4 重绘。}}
\end{frame}


\begin{frame}[t,shrink=6]{ODRP: on-demand remote paging with programmable RDMA：实验结果}
\DeckKicker{实验结果}
\SlideMeta{wang2025odrp}{E1 primary systems / Background substrate}{local\_pdf\_verified}
{\large\bfseries\color{Ink} ODRP 的实验重点是三件事同时成立：1.72x-12x utilization gain、zero MNode CPU、真实应用 0.8\%-14.6\% overhead。\par}\vspace{1.8mm}
\begin{columns}[T,totalwidth=\textwidth]
  \begin{column}{0.45\textwidth}
    \fcolorbox{Rule}{Paper}{%
  \begin{minipage}[t]{0.97\linewidth}
    \vspace{1.1mm}
    \hspace{1.4mm}\begin{minipage}{0.92\linewidth}
      \PanelTitle{讲解逻辑}
      \scriptsize \begin{itemize}
  \item 与 one-sided static 相比，ODRP 实现 1.72x 到 12x remote memory utilization 提升；与 one-sided dynamic 相比最多 1.82x。
  \item ODRP 和 one-sided static 都绕过 runtime MNode CPU；dynamic/two-sided 在多 CNode 时会饱和 MNode CPU。
  \item 真实应用 overhead 范围约 0.8\%-14.6\%，swap-intensive Kmeans 的 overhead 约 14.2\%。
  \item 8 CNodes 时 ODRP 达到 one-sided static 约 87.3\% 的 swap throughput，同时显著改善 utilization。
\end{itemize}
    \end{minipage}
    \vspace{1mm}
  \end{minipage}}
  \end{column}
  \begin{column}{0.49\textwidth}
    \fcolorbox{Rule}{Panel}{%
  \begin{minipage}[t]{0.97\linewidth}
    \vspace{1.1mm}
    \hspace{1.4mm}\begin{minipage}{0.92\linewidth}
      \PanelTitle{ODRP 关键实验数字}
      \scriptsize \textbf{utilization gain vs one-sided static}\hfill \textcolor{Accent}{\bfseries 1.72x-12x}\par\textcolor{Accent}{\rule{0.92\linewidth}{1.1mm}}\par{\tiny\textcolor{Muted}{Fig.7}}\par\vspace{1.4mm}
\textbf{runtime MNode CPU}\hfill \textcolor{Accent}{\bfseries zero}\par\textcolor{Accent}{\rule{0.04\linewidth}{1.1mm}}\par{\tiny\textcolor{Muted}{RNIC offload}}\par\vspace{1.4mm}
\textbf{application overhead}\hfill \textcolor{Accent}{\bfseries 0.8\%-14.6\%}\par\textcolor{Accent}{\rule{0.15\linewidth}{1.1mm}}\par{\tiny\textcolor{Muted}{real-world apps}}\par\vspace{1.4mm}
\textbf{Kmeans overhead}\hfill \textcolor{Accent}{\bfseries 14.2\%}\par\textcolor{Accent}{\rule{0.14\linewidth}{1.1mm}}\par{\tiny\textcolor{Muted}{swap-intensive}}\par\vspace{1.4mm}
\textbf{swap throughput at 8 CNodes}\hfill \textcolor{Accent}{\bfseries 87.3\% of static}\par\textcolor{Accent}{\rule{0.87\linewidth}{1.1mm}}\par{\tiny\textcolor{Muted}{Fig.9}}\par\vspace{1.4mm}
    \end{minipage}
    \vspace{1mm}
  \end{minipage}}
  \end{column}
\end{columns}
\vspace{1mm}
{\tiny\textcolor{Muted}{Source: 实验图按 ODRP abstract、Figure 7、Figure 8、Figure 9 和正文数值重绘。}}
\end{frame}


\begin{frame}[t,shrink=6]{ODRP: on-demand remote paging with programmable RDMA：文章评价}
\DeckKicker{文章评价}
\SlideMeta{wang2025odrp}{E1 primary systems / Background substrate}{local\_pdf\_verified}
{\large\bfseries\color{Ink} ODRP 很适合说明 SmartNIC/RNIC 可成为 memory-management actor，但它也暴露了 endpoint trust 和 state recovery 的新边界。\par}\vspace{1.8mm}
\begin{columns}[T,totalwidth=\textwidth]
  \begin{column}{0.45\textwidth}
    \fcolorbox{Rule}{Paper}{%
  \begin{minipage}[t]{0.97\linewidth}
    \vspace{1.1mm}
    \hspace{1.4mm}\begin{minipage}{0.92\linewidth}
      \PanelTitle{讲解逻辑}
      \scriptsize \begin{itemize}
  \item 优点：问题定义清晰，真正把 allocation/deallocation 从 control path 推到 data path，且用 commodity RDMA hardware。
  \item 不足：WR-chain 编程复杂，empty free queue、CNode crash、TT recovery 等 corner cases 仍需 MNode CPU 或额外 recovery 逻辑。
  \item 商业潜力：适合 remote memory、SmartNIC/DPU offload 和低 CPU memory pooling 场景。
  \item 安全边界：论文不提供 attested endpoint、encrypted link、buffer revocation policy、trusted vNIC/vSwitch 或 confidential VM ownership 证明。
\end{itemize}
    \end{minipage}
    \vspace{1mm}
  \end{minipage}}
  \end{column}
  \begin{column}{0.49\textwidth}
    \fcolorbox{Rule}{Panel}{%
  \begin{minipage}[t]{0.97\linewidth}
    \vspace{1.1mm}
    \hspace{1.4mm}\begin{minipage}{0.92\linewidth}
      \PanelTitle{落地潜力与安全边界}
      \scriptsize {\tiny
\begin{tabularx}{\linewidth}{@{}YY@{}}
\toprule
\textbf{维度} & \textbf{判断} \\
\midrule
最强贡献 & RNIC data path 同时处理 memory access 和 management \\
工程复杂度 & WR chain、EndianSwap、Modulo、recycle、crash recovery \\
可落地场景 & remote memory pool、SmartNIC/DPU offload、CPU-light MNode \\
安全缺口 & endpoint identity、link encryption、attestation、buffer lifecycle \\
\bottomrule
\end{tabularx}
}
    \end{minipage}
    \vspace{1mm}
  \end{minipage}}
  \end{column}
\end{columns}
\vspace{1mm}
{\tiny\textcolor{Muted}{Source: 评价依据 ODRP Section 4.5、Section 6、conclusion 与本仓库 confidential-I/O 证据边界。}}
\end{frame}


\begin{frame}[t,shrink=6]{Memory / I/O fabrics: CXL、PCIe IDE、RDMA：方向总结}
\DeckKicker{方向总结}
\SlideMeta{12-memory-io-fabrics}{authored direction}{local evidence synthesis}
{\large\bfseries\color{Ink} 三篇论文共同说明 data path 已经从本地 DRAM 扩展到 CXL/RDMA fabric；下一步安全问题是证明这个 fabric path 可信。\par}\vspace{1.8mm}
\begin{columns}[T,totalwidth=\textwidth]
  \begin{column}{0.45\textwidth}
    \fcolorbox{Rule}{Paper}{%
  \begin{minipage}[t]{0.97\linewidth}
    \vspace{1.1mm}
    \hspace{1.4mm}\begin{minipage}{0.92\linewidth}
      \PanelTitle{讲解逻辑}
      \scriptsize \begin{itemize}
  \item DirectCXL 优化 latency：把远端内存访问从 copy/runtime path 改成 CXL.mem address path。
  \item CXL-Tiers 优化 VM resource management：用硬件 tiering 处理 cache-line hotness，用 Memstrata 管 outlier VM。
  \item ODRP 优化 utilization 与 MNode CPU：用 RNIC WR chain 做 page-granular remote memory management。
  \item 这些都还缺 confidential I/O 组合证据：CXL IDE/PCIe IDE、SPDM/TDISP、device attestation、CCA/CoVE memory ownership 与 verifier policy。
\end{itemize}
    \end{minipage}
    \vspace{1mm}
  \end{minipage}}
  \end{column}
  \begin{column}{0.49\textwidth}
    \fcolorbox{Rule}{Panel}{%
  \begin{minipage}[t]{0.97\linewidth}
    \vspace{1.1mm}
    \hspace{1.4mm}\begin{minipage}{0.92\linewidth}
      \PanelTitle{本方向结论矩阵}
      \scriptsize {\tiny
\begin{tabularx}{\linewidth}{@{}YYY@{}}
\toprule
\textbf{论文} & \textbf{解决的 data-path 问题} & \textbf{没有解决的 confidential-I/O 问题} \\
\midrule
DirectCXL & CXL.mem latency / direct access & IDE keying、device identity、TDISP lifecycle \\
CXL-Tiers & VM tiering / performance isolation & confidential ownership 与 host placement policy 一致性 \\
ODRP & remote paging utilization / MNode CPU & attested RNIC/DPU、encrypted link、buffer revocation \\
后续组合 & CXL/RDMA fabric 可作为高性能底座 & 必须接 SPDM/TDISP/CCA/CoVE 证据链 \\
\bottomrule
\end{tabularx}
}
    \end{minipage}
    \vspace{1mm}
  \end{minipage}}
  \end{column}
\end{columns}
\vspace{1mm}
{\tiny\textcolor{Muted}{Source: 总结页使用三篇论文实验结果与 evidence ledger 的 claim-strength 边界。}}
\end{frame}
